\chapter{Les gènes et la diversité humaine}

\noindent\textit{Chaque chromosome porte, à des emplacements précis, des unités
d'information appelées gènes. Leurs différentes versions, les allèles, expliquent une
grande partie de la diversité humaine — de la couleur des yeux à des maladies
génétiques comme la drépanocytose, en passant par les groupes sanguins.}
\vspace{8pt}

\connecteBanner

\section{Le gène, unité d'information sur le chromosome}

\begin{definition}[Gène]
Un \textbf{gène} est un segment d'ADN, situé à un endroit précis d'un chromosome (son
\textbf{locus}), qui porte l'information nécessaire à un caractère héréditaire donné.
\end{definition}

\begin{definition}[Allèle]
Un \textbf{allèle} est une version particulière d'un gène. Un même gène peut exister
sous plusieurs allèles différents dans la population.
\end{definition}

\begin{propriete}[Un gène, deux exemplaires]
Chaque personne possède \textbf{deux exemplaires} de chaque gène (un hérité du père, un
hérité de la mère), portés par les deux chromosomes d'une même paire. Ces deux
exemplaires peuvent être le même allèle ou deux allèles différents.
\end{propriete}

\begin{exemple}
Le gène responsable du groupe sanguin ABO existe sous trois allèles possibles dans la
population : A, B et O.
\end{exemple}

\section{Une anomalie génique : la drépanocytose}

\begin{definition}[Drépanocytose]
La \textbf{drépanocytose} est une maladie génétique due à un allèle muté, noté
\textbf{S}, du gène qui fabrique l'hémoglobine (protéine du globule rouge qui
transporte le dioxygène). L'allèle normal est noté \textbf{A}.
\end{definition}

\begin{propriete}[Trois génotypes possibles]
\begin{itemize}
\item \textbf{AA} : les deux exemplaires du gène sont normaux $\to$ personne saine ;
\item \textbf{AS} : un exemplaire normal, un exemplaire muté $\to$ personne
\textbf{porteuse saine} (le plus souvent sans symptômes) ;
\item \textbf{SS} : les deux exemplaires sont mutés $\to$ personne \textbf{atteinte de
drépanocytose}.
\end{itemize}
\end{propriete}

\begin{propriete}[Conséquences de la maladie]
Chez une personne SS, les globules rouges prennent une forme de faucille (au lieu
d'être ronds), se bloquent dans les petits vaisseaux sanguins et se détruisent
prématurément, provoquant anémie, crises douloureuses et fatigue chronique.
\end{propriete}

\begin{illus}
\begin{tikzpicture}[scale=0.85]
\node[font=\bfseries\small] at (0,2.3) {Croisement AS $\times$ AS};
\draw[onbline] (-1.6,-1.6) rectangle (1.6,1.6);
\draw[onbline] (-1.6,0) -- (1.6,0);
\draw[onbline] (0,-1.6) -- (0,1.6);
\node[above,font=\small] at (-0.8,1.7) {A};
\node[above,font=\small] at (0.8,1.7) {S};
\node[left,font=\small] at (-1.7,0.8) {A};
\node[left,font=\small] at (-1.7,-0.8) {S};
\node[font=\small] at (-0.8,0.8) {AA};
\node[font=\small] at (0.8,0.8) {AS};
\node[font=\small] at (-0.8,-0.8) {AS};
\node[font=\small] at (0.8,-0.8) {SS};
\end{tikzpicture}
\fig{Figure — Croisement de deux parents porteurs sains (AS $\times$ AS) : 1/4 AA, 1/2
AS, 1/4 SS.}
\end{illus}

\begin{exemple}
Deux parents porteurs sains (AS) ont, à chaque grossesse, un risque de 1 sur 4 (25\,\%)
d'avoir un enfant atteint de drépanocytose (SS), 1 sur 2 (50\,\%) d'avoir un enfant
porteur sain (AS), et 1 sur 4 (25\,\%) d'avoir un enfant sain non porteur (AA).
\end{exemple}

\begin{remarque}
Piège fréquent : deux parents porteurs sains (AS) peuvent tout à fait avoir un enfant
sain (AA) ou porteur sain (AS) ; ce n'est pas parce que les deux parents portent
l'allèle S que tous leurs enfants seront malades.
\end{remarque}

\section{Les gènes et la diversité humaine : l'exemple des groupes sanguins}

\begin{propriete}[Le système ABO]
Le groupe sanguin ABO est déterminé par un gène présentant \textbf{trois allèles} : A,
B et O. Chaque personne en possède deux, ce qui donne quatre groupes phénotypiques
possibles :
\begin{itemize}
\item groupe \textbf{A} (génotype AA ou AO) ;
\item groupe \textbf{B} (génotype BB ou BO) ;
\item groupe \textbf{AB} (génotype AB) ;
\item groupe \textbf{O} (génotype OO).
\end{itemize}
\end{propriete}

\begin{propriete}[Le facteur Rhésus]
Le facteur Rhésus dépend d'un autre gène, à deux allèles : \textbf{D} (présence de
l'antigène Rhésus, Rhésus positif) et \textbf{d} (absence, Rhésus négatif). Une personne
de génotype DD ou Dd est Rhésus \textbf{positif} ; une personne dd est Rhésus
\textbf{négatif}.
\end{propriete}

\begin{exemple}
Une personne de groupe « O négatif » a le génotype OO pour le système ABO et dd pour le
système Rhésus.
\end{exemple}

\begin{remarque}
Comme pour la drépanocytose, le groupe sanguin illustre le principe général : la
diversité observable (le phénotype) résulte de la combinaison des deux allèles reçus
des parents (le génotype).
\end{remarque}

% ═══════════════════════════════════════════════════════════════
\section{L'essentiel à retenir}

\begin{synthese}
\textbf{Gène.} Segment d'ADN, situé à un locus précis, portant l'information d'un
caractère héréditaire.\\[4pt]
\textbf{Allèle.} Version particulière d'un gène ; chaque personne en possède deux
exemplaires.\\[4pt]
\textbf{Drépanocytose.} Maladie due à l'allèle muté S du gène de l'hémoglobine :
AA sain, AS porteur sain, SS malade.\\[4pt]
\textbf{Croisement AS $\times$ AS.} 1/4 AA, 1/2 AS, 1/4 SS.\\[4pt]
\textbf{Groupe sanguin ABO.} Trois allèles (A, B, O) $\to$ quatre groupes (A, B, AB,
O).\\[4pt]
\textbf{Facteur Rhésus.} Allèle D (positif) et allèle d (négatif).\\[4pt]
\textbf{Piège fréquent.} Deux parents porteurs sains (AS) peuvent avoir des enfants
sains (AA) ou porteurs (AS), pas seulement des enfants malades.
\end{synthese}

% ═══════════════════════════════════════════════════════════════
\section{Méthodes \& savoir-faire}

\methode{Interpréter un génotype}
Identifier si les deux allèles sont identiques (homozygote : AA ou SS) ou différents
(hétérozygote : AS), puis en déduire le phénotype correspondant.
\begin{exemple}
Le génotype AS correspond au phénotype « porteur sain » de la drépanocytose.
\end{exemple}

\methode{Construire un tableau de croisement}
Placer les allèles transmissibles par chaque parent en ligne et en colonne, puis
remplir chaque case avec la combinaison obtenue pour lire les proportions de
génotypes chez les enfants.
\begin{exemple}
Pour AS $\times$ AS, les quatre cases donnent AA, AS, AS, SS, soit 1/4, 1/2, 1/4.
\end{exemple}

% ═══════════════════════════════════════════════════════════════
\section{Exercices d'application}
\begin{multicols}{2}

\rubrique{Gène et allèle}
\exo{}\diff{1} Qu'est-ce qu'un gène ?

\exo{}\diff{1} Qu'est-ce qu'un allèle ?

\rubrique{La drépanocytose}
\exo{}\diff{2} Quel est l'allèle responsable de la drépanocytose ?

\exo{}\diff{2} Quel génotype correspond à une personne porteuse saine ?

\exo{}\diff{2} Quel génotype correspond à une personne malade ?

\rubrique{Groupes sanguins}
\exo{}\diff{1} Combien d'allèles possède le système ABO ?

\exo{}\diff{2} Quel génotype correspond à une personne Rhésus négatif ?

% ═══════════════════════════════════════════════════════════════
\end{multicols}

\section{Exercices d'entraînement}
\begin{multicols}{2}

\rubrique{La drépanocytose}
\exo{}\diff{2} Deux parents AS ont trois enfants. Quelle est, pour chaque enfant, la
probabilité qu'il soit atteint de drépanocytose ?

\exo{}\diff{3} Un couple AA $\times$ AS peut-il avoir un enfant SS ? Justifier à l'aide
d'un tableau de croisement.

\exo{}\diff{3} Expliquer pourquoi deux parents porteurs sains (AS) peuvent avoir un
enfant en parfaite santé et non porteur (AA).

\rubrique{Groupes sanguins}
\exo{}\diff{2} Une personne de groupe A peut avoir le génotype AA ou AO. Expliquer
pourquoi deux génotypes différents donnent le même groupe sanguin.

\exo{}\diff{3} Un enfant est de groupe O. Ses deux parents peuvent-ils être de groupe A
? Justifier à l'aide des génotypes possibles.

% ═══════════════════════════════════════════════════════════════
\end{multicols}

\section{Sujets type BEPC}

\sujetbac[Gènes et drépanocytose]
\begin{enumerate}[label=\arabic*)]
\item Définis les mots « gène » et « allèle ».
\item Cite les trois génotypes possibles pour la drépanocytose et le phénotype
correspondant à chacun.
\item Deux parents porteurs sains (AS) attendent un enfant. Construis le tableau de
croisement et donne les proportions de génotypes attendues.
\item Explique pourquoi un enfant AA peut naître de deux parents porteurs sains.
\end{enumerate}

\sujetbac[Diversité génétique : les groupes sanguins]
\begin{enumerate}[label=\arabic*)]
\item Cite les trois allèles du système ABO.
\item Cite les quatre groupes sanguins possibles.
\item Donne les deux allèles du système Rhésus et le phénotype associé à chacun.
\item En quoi les groupes sanguins illustrent-ils la diversité génétique humaine ?
\end{enumerate}

% ═══════════════════════════════════════════════════════════════
\section{Corrigés}
\begin{multicols}{2}

\corrige{de l'exercice 1}
Un segment d'ADN, situé à un endroit précis d'un chromosome, qui porte l'information
d'un caractère héréditaire.

\corrige{de l'exercice 2}
Une version particulière d'un gène.

\corrige{de l'exercice 3}
L'allèle S.

\corrige{de l'exercice 4}
Le génotype AS.

\corrige{de l'exercice 5}
Le génotype SS.

\corrige{de l'exercice 6}
Trois allèles : A, B et O.

\corrige{de l'exercice 7}
Le génotype dd.

\corrige{de l'exercice 8}
Chaque enfant a, indépendamment des autres, une probabilité de 1/4 (25\,\%) d'être
atteint de drépanocytose.

\corrige{de l'exercice 9}
Non : le parent AA ne peut transmettre que l'allèle A ; le tableau de croisement donne
uniquement des enfants AA ou AS, jamais SS.

\corrige{de l'exercice 10}
Chaque parent porteur sain transmet soit l'allèle A soit l'allèle S à son enfant ; si
les deux parents transmettent chacun leur allèle A, l'enfant reçoit deux allèles A et
est donc AA, sain et non porteur.

\corrige{de l'exercice 11}
Parce que l'allèle O est récessif : chez un génotype AO, l'allèle A s'exprime seul et
détermine le groupe A, tout comme chez un génotype AA.

\corrige{de l'exercice 12}
Oui : un parent de groupe A peut avoir le génotype AO (et non AA) ; si les deux parents
sont AO, ils peuvent chacun transmettre l'allèle O et donner un enfant de génotype OO,
donc de groupe O.

\corrige{du sujet 1}
1) Un gène est un segment d'ADN portant l'information d'un caractère héréditaire ; un
allèle est une version particulière de ce gène.\\
2) AA (sain), AS (porteur sain), SS (malade).\\
3) Le tableau donne AA, AS, AS, SS, soit 1/4 AA, 1/2 AS, 1/4 SS.\\
4) Parce que chaque parent porteur sain peut transmettre l'allèle A ; si les deux
parents transmettent leur allèle A, l'enfant est AA.

\corrige{du sujet 2}
1) A, B et O.\\
2) A, B, AB et O.\\
3) L'allèle D (Rhésus positif) et l'allèle d (Rhésus négatif).\\
4) Parce que la combinaison de plusieurs allèles possibles pour un même gène produit
plusieurs phénotypes différents dans la population, comme pour de nombreux autres
caractères héréditaires.

\end{multicols}
\exercicesBanner
